\section{Step 26 Finish the vector column round}\label{chacha-s26}

Apply the remaining two groups with left rotations eight and seven. The row layout still represents ordinary state order. Each lane has now completed the same quarter-round as the corresponding scalar column call.

\subsection{Write and inspect}
Replace the body of region \texttt{msa\_column\_b}. Retain its \texttt{\# BEGIN} and \texttt{\# END} comments and the supplied labels outside them.

\Needspace{15\baselineskip}
\MARSCodeLabel{MIPS assembly}
\begin{minted}[samepage=true]{mars-idea}
    addv.w $w0, $w0, $w1
    xor.v $w3, $w3, $w0
    slli.w $w8, $w3, 8
    srli.w $w9, $w3, 24
    or.v $w3, $w8, $w9
    addv.w $w2, $w2, $w3
    xor.v $w1, $w1, $w2
    slli.w $w8, $w1, 7
    srli.w $w9, $w1, 25
    or.v $w1, $w8, $w9
\end{minted}

Use recipe \texttt{s26-1} with the cumulative program through this edit. Inspect its named checkpoint before continuing to the temporary or normal return.

\subsection{Expected observation}
The four working rows equal the scalar first column-round state.

\Needspace{9\baselineskip}
\subsection{Check your understanding}
What do these vector groups complete?

\begin{itemize}
\item \textbf{A} Twenty rounds of one block.
\item \textbf{B} Four encrypted messages.
\item \textbf{C} Four quarter-rounds, one in each lane.
\end{itemize}
